\documentclass[11pt]{article}

\usepackage[T1]{fontenc}
\usepackage[utf8]{inputenc}
\usepackage{lmodern}
\usepackage{amsmath}
\usepackage{amssymb}
\usepackage{booktabs}
\usepackage{array}
\usepackage{geometry}
\usepackage{hyperref}
\usepackage{url}
\usepackage{microtype}

\geometry{letterpaper, margin=1in}
\hypersetup{
  colorlinks=true,
  linkcolor=blue,
  citecolor=blue,
  urlcolor=blue,
  pdftitle={Aliasing-Aware RTNeural-Compatible WaveNet Modeling of Guitar Amplifier and Pedal Captures},
  pdfauthor={Shortwav Labs}
}

\newcommand{\esr}{\mathrm{ESR}}
\newcommand{\asr}{\mathrm{ASR}}
\newcommand{\rms}{\mathrm{RMS}}
\newcommand{\rtf}{\mathrm{RTF}}

\title{Aliasing-Aware RTNeural-Compatible WaveNet Modeling of Guitar Amplifier and Pedal Captures}

\author{
Shortwav Labs\\
\small RTNeural Trainer Project\\
\small \url{https://github.com/shortwavlabs/rtneural-trainer}
}

\date{June 29, 2026}

\begin{document}
\maketitle

\begin{abstract}
We report an applied study of neural guitar amplifier and pedal modeling under
a strict real-time export constraint: trained models must be representable as
RTNeural JSON, validate against TensorFlow/Keras reference inference, and run as
causal low-latency DSP in a native C++ audio runtime. The study began with
dense, recurrent, and convolutional candidates, but the strongest end-to-end
evidence in our internal 48 kHz mono paired-capture dataset favored
WaveNet-style temporal convolutional networks (TCNs). We evaluate clean,
edge-of-breakup, crunch, rhythm, lead, and overdrive-pedal captures using
waveform error, residual statistics, native real-time factor, and a
probe-based aliasing-to-signal ratio (ASR). The best high-gain result came from
an A2-inspired sequential TCN using mixed kernel sizes, non-power-of-two
dilations, and PReLU nonlinearities while remaining compatible with current
RTNeural dynamic JSON. On a difficult high-gain rhythm capture, this model
improved preview error-to-signal ratio (ESR) from 0.0646 for the previous
smoothed-tanh quality model to 0.0381 after continuation, while reducing
probe-measured average ASR from 0.0419 to 0.0145 and benchmarking at a
worst-case native Eigen real-time factor of approximately 6.62 on the
reference system. Real hardware captures later produced stronger export
packages: clean/edge, overdrive pedal, and rhythm exports reached ESR values of
0.00217, 0.00043, and 0.00355 respectively, all passing native RTNeural
validation with low probe ASR. The results support a practical, scoped
conclusion: under current RTNeural JSON and product-workflow constraints,
WaveNet-family causal TCNs are the best-tested architecture family for this
trainer, while exact NAM A2-style residual/skip topology remains future
runtime work.
\end{abstract}

\section{Introduction}

Neural modeling of guitar amplifiers and pedals is a system-identification
problem with unusually unforgiving perceptual constraints. A fixed amp or pedal
setting maps a dry direct-input signal \(x[n]\) to a processed target
\(y[n]\), but the mapping is nonlinear, level-dependent, often highly
compressive, and sensitive to latency errors of only a few samples. High-gain
tones further expose residual energy in the upper bands, where small spectral
mistakes can become audible as fizz, foldback, or loss of pick definition.

The goal of RTNeural Trainer is not merely to minimize a training loss. The
model must be usable as real-time DSP: it must be causal, streamable,
exportable to RTNeural, parity-checked against the training backend,
benchmarked in a native runtime, and loadable in a DAW plugin. This engineering
constraint rules out treating the training graph as an arbitrary offline model.
It also makes runtime validation a first-class part of model evaluation.

This paper summarizes the DSP research performed while building the trainer.
Our main contributions are:

\begin{itemize}
  \item an end-to-end RTNeural-compatible training and export workflow for
  paired guitar amp/pedal captures;
  \item a comparison of product-facing WaveNet-style TCN presets across clean,
  edge, crunch, rhythm, lead, and overdrive-pedal cases;
  \item an aliasing-to-signal ratio diagnostic for exported RTNeural JSON
  models;
  \item an A2-inspired sequential TCN that improves difficult high-gain
  captures without requiring new RTNeural graph topology; and
  \item a set of capture and validation lessons for real-time neural audio
  effects.
\end{itemize}

The central claim is deliberately scoped. We do not prove that TCNs dominate
all recurrent or convolutional models for all amplifier modeling tasks. We show
that, in this internal dataset and under the current RTNeural JSON export
constraint, WaveNet-family causal TCNs produced the strongest end-to-end
product evidence.

\section{Related Work}

WaveNet introduced dilated causal convolutional stacks for raw audio modeling
\cite{oord2016wavenet}. Subsequent neural audio-effects and amplifier-modeling
work has used recurrent models, convolutional models, spectral losses,
pre-emphasis, and finite-memory feed-forward architectures for real-time
nonlinear system approximation \cite{damskagg2019tube,steinmetz2022auraloss}.
PANAMA frames amplifier modeling as a parametric and active-learning problem,
where control settings are sampled strategically and the model is conditioned
on those settings \cite{grotschla2025panama}. That broader parametric
setting remains future work here; the present trainer models one fixed capture
setting per project.

Neural Amp Modeler (NAM) is an important practical reference for this work
\cite{namrepo,namplugin,namcore}. NAM separates training/export from real-time
plugin inference and treats the exported model package as a runtime contract.
Recent NAM A2 models use WaveNet-like residual and skip structures with
specialized C++ runtime paths. Current RTNeural dynamic JSON can host
sequential Conv1D and PReLU stacks \cite{rtneural}, but it does not directly
express A2-style residual fan-out, skip accumulation, layer-array heads, or
slimmable submodel selection. This motivated an A2-inspired sequential preset
rather than direct NAM conversion.

Aliasing is another relevant thread. Nonlinear models generate harmonics that
can fold above Nyquist, and high-gain guitar models stress this behavior. Work
on smoothing neural activations for amp modeling suggests that activation shape
can trade waveform accuracy against aliasing behavior
\cite{steinmetz2025aliasing}. We therefore added explicit ASR diagnostics and
smoothed-tanh WaveNet variants, but treat ASR as a warning signal rather than a
psychoacoustic score.

\section{Problem Formulation}

Let \(x[n]\) be the prepared dry input and \(y[n]\) the prepared target. After
channel policy, resampling policy, polarity handling, and latency alignment,
both are mono 48 kHz floating-point sequences. A causal finite-memory model
produces
\begin{equation}
  \hat{y}[n] = f_\theta(x[n], x[n-1], \ldots, x[n-R+1]),
\end{equation}
where \(R\) is the effective receptive field. For a sequential stack of causal
Conv1D layers with kernel sizes \(k_\ell\), dilations \(d_\ell\), and no
striding, the receptive field is
\begin{equation}
  R = 1 + \sum_\ell (k_\ell - 1)d_\ell .
\end{equation}

The residual is
\begin{equation}
  e[n] = y[n] - \hat{y}[n].
\end{equation}
The primary waveform metric is error-to-signal ratio:
\begin{equation}
  \esr = \frac{\sum_n e[n]^2}{\max(\sum_n y[n]^2, \epsilon)} .
\end{equation}
We also report mean absolute error, root mean square error, peak residual, and
correlation. Prediction level is tracked by
\begin{equation}
  \rho_{\rms} = \frac{\rms(\hat{y})}{\max(\rms(y), \epsilon)} ,
\end{equation}
which guards against underpowered predictions that appear deceptively stable.

Three metric spans appear in our reports. Window validation is computed on
held-out training windows. Stream validation is computed on a longer
continuous segment to catch state and continuity failures. Preview or
state-continuous metrics are computed on rendered preview WAVs from a saved
checkpoint or export and are the main user-facing comparison metrics in this
paper unless otherwise noted.

\section{Methods}

\subsection{Capture preparation}

The capture contract is a paired dry/processed recording from the same
performance. Captures are prepared as 48 kHz mono WAVs, with float32 preserved
when available. Stereo or dual-mono sources are mixed to mono for the current
trainer. The preparation stage reports clipping, duration, level, channel
policy, latency candidates, confidence, and polarity.

Latency is treated as an experimental variable. If the target is shifted by
\(L\) samples, then the prepared pair is
\begin{equation}
  x_{\mathrm{aligned}}[n] = x[n], \qquad
  y_{\mathrm{aligned}}[n] = y[n+L].
\end{equation}
For high-gain captures the processed waveform can correlate weakly with the
dry DI, producing several plausible offsets separated by only a few samples.
The application therefore exposes top candidate offsets and allows manual
sample nudging. In practice we found that a short transient preamble at the
start of the DI---palm mutes, single-note plucks, and hard/soft attacks---is
valuable for alignment.

\subsection{Training protocol}

Training uses TensorFlow/Keras only. The default optimizer is Adam, with a
default seed of 1337 unless overridden. Default batch size is 16 and default
sequence length is 1024 samples, though product recipes can override these
values. Candidate windows are sampled across the prepared file; validation
windows are fixed, while training windows may be resampled between epochs.

Checkpoint selection uses a composite validation score:
\begin{equation}
  s = \esr_{\mathrm{stream}} + 0.25\,\esr_{\mathrm{window}} + p_{\mathrm{under}},
\end{equation}
where \(p_{\mathrm{under}}\) penalizes severely underpowered predictions. Early
stopping monitors this score with default patience 5 and minimum improvement
\(10^{-4}\). The learning rate scheduler uses ReduceLROnPlateau with factor
0.5, a minimum learning rate of \(10^{-6}\), and patience derived from the
early-stopping patience.

Losses include waveform ESR/MSE variants, pre-emphasis MSE, and a
multi-resolution STFT plus pre-emphasis loss. The pre-emphasis coefficient is
0.95, with pre-emphasis loss weight 0.35. The multi-resolution STFT frame sizes
are 256, 1024, and 2048 samples, with STFT weight 0.02 and log-magnitude weight
0.05. These values are implementation settings rather than claimed optima.

\subsection{Architectures}

The project initially included dense, GRU, LSTM, shallow Conv1D,
BatchNorm/PReLU, and hybrid presets. Those paths remain useful as internal
export fixtures, but product-facing recipes have moved to WaveNet-family TCNs.
The main presets are summarized in Table~\ref{tab:presets}.

\begin{table}[t]
\centering
\small
\begin{tabular}{@{}lrrp{2.2cm}p{3.0cm}p{2.2cm}@{}}
\toprule
Preset & Layers & Filters & Kernels & Dilations & Activation \\
\midrule
Fast & 6 & 12 & 3 & \(1,\ldots,32\) powers of two & tanh \\
Balanced & 8 & 16 & 3 & \(1,\ldots,128\) powers of two & tanh \\
Quality & 10 & 20 & 3 & \(1,\ldots,512\) powers of two & tanh \\
Quality Tanh 1.5 & 10 & 20 & 3 & \(1,\ldots,512\) powers of two & \(\tanh(x/1.5)\) \\
Clean & 10 & 8 & 7 & \(1,\ldots,512\) powers of two & linear hidden \\
Edge & 10 & 8 & 7 & \(1,\ldots,512\) powers of two & \(\tanh(x/1.8)\) \\
A2 PReLU & 12 & 16 & 6, 15 mixed & non-power-of-two mixed & PReLU \\
\bottomrule
\end{tabular}
\caption{Principal RTNeural-compatible TCN presets. The shorthand
\(1,\ldots,512\) denotes powers-of-two dilations. The A2 PReLU row uses the
mixed sequence \(1,3,7,17,41,101,239,1,3,7,17,41\).}
\label{tab:presets}
\end{table}

These models are intentionally sequential:
\begin{equation}
  x \rightarrow [\mathrm{causal\ Conv1D} \rightarrow
  \mathrm{activation}]^N \rightarrow \mathrm{output}.
\end{equation}
They are not full residual/skip WaveNet graphs. That limitation preserves
RTNeural dynamic-JSON compatibility, but it also limits how much depth can be
added before optimization becomes difficult. A failed high-gain experiment with
one additional plain tanh dilation stage suggested that larger receptive field
alone does not solve difficult rhythm captures.

\subsection{Aliasing-to-signal ratio}

For each exported RTNeural JSON model, we render deterministic sine probes at
requested frequencies 1250, 2500, and 5000 Hz. The input amplitude is 0.5, with
2048 warmup samples and a 4096-sample analysis window. The actual frequency is
snapped to the nearest FFT bin so the fundamental is on-bin. After subtracting
the mean, the analyzer computes FFT power and identifies harmonic bins. It then
defines
\begin{align}
  E_{\mathrm{harm}} &= \sum_{b\in \mathcal{H}} P[b],\\
  E_{\mathrm{alias}} &= \sum_{b>0} P[b] - E_{\mathrm{harm}},\\
  \asr &= \frac{E_{\mathrm{alias}}}{\max(E_{\mathrm{harm}}, \epsilon)} .
\end{align}
Worst-case ASR below 0.02 is treated as low aliasing in the current probes;
values below 0.08 produce a review warning; larger values produce a stronger
warning. These thresholds are not audibility thresholds. They are engineering
diagnostics to focus listening tests.

\subsection{Native validation and plugin smoke testing}

Exports are gated by Keras-to-RTNeural JSON parity, native RTNeural validation,
and native benchmark results. The validator tests block sizes 16, 32, 64, 128,
256, and 512 samples, mono and stereo when applicable, and records worst-case
real-time factor:
\begin{equation}
  \rtf = \frac{\mathrm{audio\ duration\ processed}}
              {\mathrm{wall\ time\ elapsed}} .
\end{equation}
An RTF above one is faster than real time, but product headroom requires a
larger margin because DAWs may run multiple instances and small buffers.

A JUCE Audio Unit debug plugin was used for smoke testing \cite{juce}. It
loads exported RTNeural package folders, stages model objects off the audio
thread, and keeps \texttt{processBlock()} free of parsing, file I/O, locks, and
heap allocation. These plugin tests were performed on a MacBook Pro M5 Max
reference machine and should not be interpreted as broad release qualification.

\section{Experiments}

The main internal capture families are listed in Table~\ref{tab:captures}. All
audio was prepared at 48 kHz mono for training.

\begin{table}[t]
\centering
\small
\begin{tabular}{@{}lp{3.0cm}p{6.6cm}@{}}
\toprule
Family & Duration & Notes \\
\midrule
DI2 & 151.6 s & Same dry DI rendered to clean, crunch, rhythm, edge, lead,
and overdrive-pedal targets. \\
DI3/CLEAN3 & 613.08 s & Long clean capture with known 13-sample DAW latency. \\
DI3-B/RHYTHM3-B & 444.25 s & Trimmed high-gain rhythm capture; 10-sample
estimated latency with low confidence. \\
DI4/RHYTHM4 & 158.8 s & Short high-gain rhythm control capture; 9-sample
estimated latency. \\
Hardware trio & focused captures & Real amp clean/edge, real overdrive pedal,
and real amp rhythm exports. \\
\bottomrule
\end{tabular}
\caption{Capture families used in the internal experiments.}
\label{tab:captures}
\end{table}

The DI2 family established tone-dependent difficulty. WaveNet quality achieved
preview ESR 0.0052 on clean, 0.0115 on crunch, 0.1100 on rhythm, 0.0045 on
edge, and 0.0021 on the overdrive pedal. Lead was the primary outlier:
balanced slightly beat quality at ESR 0.0810 versus 0.0821. Shallow stacked
Conv/PReLU baselines were fast but trailed WaveNet-family presets across the
reported product comparisons.

\section{Results}

\subsection{High-gain RHYTHM4 progression}

Table~\ref{tab:rhythm4} shows the most informative high-gain progression. The
plain balanced model plateaued. The quality model fit the waveform much better
but produced a strong ASR warning. The smoothed-tanh quality model improved
both ESR and probe ASR. The A2-inspired PReLU model then improved both metrics
again, at a runtime and model-size cost.

\begin{table}[t]
\centering
\small
\begin{tabular}{@{}lrrrrrr@{}}
\toprule
Preset & ESR & RMSE & Corr. & Worst ASR & Avg. ASR & Native RTF \\
\midrule
Balanced & 0.6309 & 0.0740 & 0.6109 & -- & -- & -- \\
Quality & 0.0713 & 0.0249 & 0.9640 & 0.290 & 0.100 & 12.17 \\
Quality Tanh 1.5 & 0.0646 & 0.0237 & 0.9674 & 0.0670 & 0.0419 & 11.74 \\
A2 PReLU & 0.0440 & 0.0196 & 0.9778 & 0.0354 & 0.0205 & 6.54 \\
A2 PReLU continued & 0.0381 & 0.0182 & 0.9808 & 0.0201 & 0.0145 & 6.62 \\
\bottomrule
\end{tabular}
\caption{RHYTHM4 high-gain progression. Missing ASR/RTF values indicate runs
not used as export candidates in the benchmarked export study.}
\label{tab:rhythm4}
\end{table}

This sequence does not isolate which A2-inspired change caused the
improvement. Kernel sizes, dilation pattern, activation type, training dynamics,
and continuation all differ. The result should therefore be read as evidence
for the combined preset, not as a controlled ablation of individual features.

\subsection{Clean and edge-of-breakup behavior}

The real clean amp project exposed an architecture mismatch. A2 PReLU initially
improved but then diverged from full-stream validation, despite stable
130-sample alignment and polarity handling. A clean long-field linear preset
improved the run to ESR 0.0675 and correlation 0.9663. The capture was
technically clean for single coils but behaved closer to edge-of-breakup with
humbuckers; adding gentle hidden nonlinearity in the edge preset reduced ESR to
0.00362 and correlation to 0.9982. A larger edge-detail variant did not
improve this result, suggesting that the key correction was light nonlinearity
rather than more capacity.

\subsection{Real hardware exports}

The hardware export trio is summarized in Table~\ref{tab:hardware}. All three
exports passed native RTNeural validation and reported low probe ASR.

\begin{table}[t]
\centering
\small
\begin{tabular}{@{}llrrrr@{}}
\toprule
Export & Preset & ESR & Worst ASR & Avg. ASR & Native RTF \\
\midrule
Clean/edge amp & Edge & 0.00217 & 0.00673 & 0.00251 & 20.72 \\
Overdrive pedal & A2 PReLU & 0.00043 & 0.00195 & 0.00080 & 6.49 \\
Rhythm amp & A2 PReLU & 0.00355 & 0.00913 & 0.00359 & 6.63 \\
\bottomrule
\end{tabular}
\caption{Real hardware export results. Native RTF values are reference-system
Eigen-backend results.}
\label{tab:hardware}
\end{table}

These hardware captures trained more reliably than several earlier DAW
amp-simulation renders. Plausible contributors include the matured capture
workflow, float32 preparation, shorter focused material, consistent reamping,
and the analog capture path's natural band limiting. These explanations remain
hypotheses; matched experiments would be needed to isolate causes.

\section{Discussion}

\subsection{Why TCNs fit the product constraint}

The successful presets combine three properties: finite memory, causal
streaming, and RTNeural compatibility. Dilation expands temporal context
without sample-by-sample recurrence, and Conv1D exports naturally to current
RTNeural JSON. Smaller Conv models were fast but left residual structure in
important tones. Recurrent models were considered during development, but the
strongest validated export evidence in this workflow came from TCNs.

The models in this paper are not full WaveNet graphs with residual and skip
topology. They are sequential TCNs shaped by export constraints. That
constraint is useful: it forces every architecture to survive the same native
validation and real-time benchmark path.

\subsection{Latency as experimental control}

Latency is a major bottleneck. A few samples of error can dominate
high-frequency residuals, and high-gain distortion weakens simple correlation
between dry and processed signals. The trainer therefore treats latency as an
interactive measurement problem, not a hidden preprocessing step. Future work
should combine transient alignment, envelope correlation, band-limited scoring,
phase methods, and post-training residual-shift diagnostics.

\subsection{ESR, ASR, and listening}

ESR is necessary but insufficient. A low-ESR model can still alias, mismatch
level, or produce unacceptable residual spectra. ASR provides a deterministic
probe-specific stress test for nonlinear foldback, but it is not a perceptual
model. The appropriate export gate is multi-objective: backend parity, native
validation, runtime headroom, preview metrics, probe ASR, residual listening,
and user judgment.

\subsection{Runtime scope}

The A2 PReLU result ran comfortably on the reference machine and in a JUCE
debug plugin smoke test, including a four-instance Logic Pro test at a
32-sample buffer. This is encouraging but not production qualification. The
next runtime work is controlled plugin benchmarking across weaker systems,
release builds, sample rates, block sizes, plugin validation tools
\cite{pluginval}, and full pedal-plus-amp-plus-IR chains. Static RTNeural
\texttt{ModelT} specializations or fused residual TCN layers may eventually be
necessary.

\section{Limitations}

This study is an internal applied research report rather than a fully
controlled public benchmark. Most comparisons are single-seed product
experiments. Dataset diversity is limited to the available captures and
hardware. ASR thresholds are engineering heuristics rather than calibrated
audibility thresholds. Some high-gain captures have low latency confidence.
Runtime evidence is strongest on Apple Silicon with the Eigen backend. Finally,
although the A2-inspired sequential preset is effective, it is not equivalent
to NAM A2 topology.

\section{Future Work}

The most immediate research tasks are repeated-seed comparisons for the best
WaveNet variants, matched latency-offset ablations, listening-calibrated ASR
thresholds, and controlled plugin benchmarks. A day-2 RTNeural runtime
extension should add a fused \texttt{a2\_wavenet} or \texttt{residual\_tcn}
layer with preallocated state, residual/skip accumulation, and constrained JSON
schema. Longer-term work includes NAM-teacher distillation, plugin-side
oversampling or higher-rate model support, and PANAMA-style parametric active
learning for controllable amplifier settings.

\section{Conclusion}

Under the current 48 kHz mono paired-capture workflow and RTNeural JSON export
constraint, WaveNet-family causal TCNs are the best-tested architecture family
for RTNeural Trainer. The strongest high-gain evidence is an A2-inspired
sequential PReLU TCN that substantially improves RHYTHM4 waveform fit and
probe ASR while remaining compatible with current RTNeural exports. The
strongest practical evidence is the real hardware clean/drive/rhythm export
trio, which passed native validation with low probe ASR and reference-system
runtime headroom.

The broader lesson is that neural amp modeling is not just a training problem.
Capture discipline, latency alignment, spectral loss, aliasing diagnostics,
native validation, and DAW runtime behavior all determine whether a model is
musically useful.

\section*{Data and Code Availability}

The trainer implementation, native validator, documentation, and experimental
notes are maintained at
\url{https://github.com/shortwavlabs/rtneural-trainer}. Some source captures
used in this study are not included in the repository because they contain
performance audio and third-party tone references.

\section*{Acknowledgments}

This work builds on the RTNeural, Neural Amp Modeler, JUCE, and audio-ML
research communities. We thank the maintainers of those projects for making
their code and documentation available.

\begin{thebibliography}{99}

\bibitem{oord2016wavenet}
A.~van~den Oord, S.~Dieleman, H.~Zen, K.~Simonyan, O.~Vinyals, A.~Graves,
N.~Kalchbrenner, A.~Senior, and K.~Kavukcuoglu.
\newblock WaveNet: A generative model for raw audio.
\newblock \emph{arXiv preprint arXiv:1609.03499}, 2016.

\bibitem{damskagg2019tube}
E.-P. Damsk\"agg, L.~Juvela, E.~Thuillier, and V.~V\"alim\"aki.
\newblock Deep learning for tube amplifier emulation.
\newblock \emph{ICASSP}, 2019. arXiv:1811.00334.

\bibitem{steinmetz2022auraloss}
C.~J. Steinmetz and J.~D. Reiss.
\newblock auraloss: Audio-focused loss functions in PyTorch.
\newblock In \emph{Digital Music Research Network One-day Workshop}, 2020.

\bibitem{grotschla2025panama}
F.~Gr\"otschla, L.~A. Lanzend\"orfer, L.~Jiao, and R.~Wattenhofer.
\newblock Parametric neural amp modeling with active learning.
\newblock \emph{arXiv preprint arXiv:2507.02109}, 2025.

\bibitem{steinmetz2025aliasing}
R.~Sato and J.~O. Smith III.
\newblock Aliasing reduction in neural amp modeling by smoothing activations.
\newblock \emph{arXiv preprint arXiv:2505.04082}, 2025.

\bibitem{namrepo}
S.~Atkinson.
\newblock Neural Amp Modeler.
\newblock \url{https://github.com/sdatkinson/neural-amp-modeler}.

\bibitem{namplugin}
S.~Atkinson.
\newblock NeuralAmpModelerPlugin.
\newblock \url{https://github.com/sdatkinson/NeuralAmpModelerPlugin}.

\bibitem{namcore}
S.~Atkinson.
\newblock NeuralAmpModelerCore.
\newblock \url{https://github.com/sdatkinson/NeuralAmpModelerCore}.

\bibitem{rtneural}
J.~Chowdhury.
\newblock RTNeural.
\newblock \url{https://github.com/jatinchowdhury18/RTNeural}.

\bibitem{juce}
Raw Material Software.
\newblock JUCE.
\newblock \url{https://juce.com/}.

\bibitem{pluginval}
Tracktion.
\newblock pluginval.
\newblock \url{https://github.com/Tracktion/pluginval}.

\end{thebibliography}

\end{document}
